\documentclass[10pt]{article}

\usepackage[margin=0.75in]{geometry}
\usepackage{amsmath,amsthm,amssymb,bm}
\usepackage{xcolor}
\usepackage{cancel}
\usepackage{graphicx}
\usepackage{changepage}
\usepackage{circuitikz}
\usepackage{pgfplots}
\usepackage{minted}
\usepackage{physics}
\usepackage{tikz}
\usepackage{tikz-3dplot}
\usepackage{hyperref}
\usepackage{cleveref}
\usepackage{siunitx}
\usepackage{fontspec}
\usepackage{relsize}
\usepackage{subfig}
\usepackage{todonotes}
\usepackage{contour}
\usepackage{multicol, multirow, booktabs}
\usepackage[breakable]{tcolorbox}
\usepackage[inline]{enumitem}

\theoremstyle{definition}
\newtheorem{problem}{Problem}
\newtheorem{soln}{Solution}

\pgfplotsset{compat=newest}
\usetikzlibrary{lindenmayersystems}
\usetikzlibrary{arrows}
\usetikzlibrary{calc}
\usetikzlibrary{positioning, fit}
\usetikzlibrary{3d, perspective}
\usetikzlibrary{math}
\usetikzlibrary{fadings}
\usetikzlibrary{angles,quotes} 
\usetikzlibrary{decorations.markings}

\definecolor{incolor}{HTML}{303F9F}
\definecolor{outcolor}{HTML}{D84315}
\definecolor{cellborder}{HTML}{CFCFCF}
\definecolor{cellbackground}{HTML}{F7F7F7}
\newcommand{\ui}{\hat{i}}
\newcommand{\uj}{\hat{j}}
\newcommand{\uk}{\hat{k}}
\newcommand{\ux}{\hat{\mathbf{x}}}
\newcommand{\uy}{\hat{\mathbf{y}}}
\newcommand{\uz}{\hat{\mathbf{z}}}
\newcommand{\uv}[1]{\hat{\bv{{#1}}}}
\newcommand{\pr}[1]{{#1^\prime}}
\newcommand{\pd}[2][]{{\frac{\partial^{#1}}{\partial {{#2}^{#1}}}}}
\newcommand{\dif}[2][]{{\frac{d^{#1}}{d{{#2}^{#1}}}}}
\newcommand{\grd}{\nabla}

\pgfdeclarelayer{background}  
\pgfsetlayers{background,main}
\AtBeginDocument{\RenewCommandCopy\qty\SI}

\makeatletter
\newcommand{\boxspacing}{\kern\kvtcb@left@rule\kern\kvtcb@boxsep}
\makeatother
\newcommand{\prompt}[4]{
    \ttfamily\llap{{\color{#2}[#3]:\hspace{3pt}#4}}\vspace{-\baselineskip}
}

\newcommand{\thevenin}[2]{
  \begin{center}
    \begin{circuitikz} \draw
      (0,0) -- (2,0) to[battery1, l_=$V_{Th}\eq#1$] (2,2) 
      to[resistor, l_=$R_{Th}\eq#2$] (0,2)
      ;
      \draw [o-] (-.07,2.079);
      \draw [o-] (-.07,0.079);
    \end{circuitikz}
  \end{center}
}

\newcommand{\norton}[2]{
  \begin{center}
    \begin{circuitikz} \draw
      (0,0) -- (3,0) to[american current source, l_=$I_{N}\eq#1$] (3,2) -- (0,2) (2,0)
      to[resistor, l=$R_{N}\eq#2$] (2,2)
      ;
      \draw [o-] (-.07,2.079);
      \draw [o-] (-.07,0.079);
    \end{circuitikz}
  \end{center}
}

\DeclareMathOperator{\Div}{div}
\DeclareMathOperator{\Curl}{curl}
\DeclareMathOperator{\sgn}{sgn}

\newcommand{\highlight}[1]{\colorbox{yellow}{$\displaystyle #1$}}
\newcommand{\ti}[1]{\widetilde{#1}}
\renewcommand{\div}[1]{\bv{\nabla}\cdot #1}
\renewcommand{\curl}[1]{\bv{\nabla}\times #1}

\newfontface{\Kaufmann}{Kaufmann}
\DeclareTextFontCommand{\kf}{\Kaufmann}
\newcommand{\scriptr}{\fontsize{12pt}{12pt}\kf{r}\fontsize{10pt}{10pt}}

\newfontface{\KaufmannB}{Kaufmann Bold}
\DeclareTextFontCommand{\kfb}{\KaufmannB}
\newcommand{\bscriptr}{\fontsize{12pt}{12pt}\kfb{r}\fontsize{10pt}{10pt}}
\newcommand{\justif}[2]{&{#1}&\text{#2}}
\newcommand{\bv}[1]{\bm{\mathrm{#1}}}

\title{Physics 4220H: Assignment III}
\author{Jeremy Favro (0805980) \\ Trent University, Peterborough, ON, Canada}
\date{\today}

\begin{document}
\maketitle

% PROBLEM 1
\begin{problem}
Starting from the following expression for (the square of) the complex
propagation constant in a conductor,
$$
  \ti{k}^2=\mu\epsilon\omega^2+i\mu\sigma\omega=\mu\omega\left[\epsilon\omega+i\sigma\right]
$$
Go through the algebra to find expressions for the real ($\mathcal{k}$) and imaginary ($\kappa$) parts
of the propagation constant. i.e. find $\mathcal{k}$ and $\kappa$, where
$\ti{k}=\mathcal{k}+\kappa$.
\end{problem}
\begin{soln}
  To do this I will write $\ti{k}$ in polar form, apply the square root, then use
  the half-angle formulae
  $$\cos(x/2)=\sqrt{\frac{1+\cos(x)}{2}}\quad\sin(x/2)=\sqrt{\frac{1-\cos(x)}{2}}$$
  and the formulae (derivable by considering a triangle with sides $1$, $x$ and $\sqrt{1+x^2}$)
  $$\cos(\arctan(x))=\frac{1}{\sqrt{1+x^2}}\quad\sin(\arctan(x))=\frac{x}{\sqrt{1+x^2}}.$$
  To start, we can write
  $$\ti{k}^2=Ke^{i\phi}$$
  with
  $$K=\abs{\ti{k}}=\mu\omega\sqrt{(\epsilon\omega)^2+\sigma^2}\quad \phi=\arctan(\frac{\sigma}{\epsilon \omega}).$$
  We then write
  \begin{align*}
    \ti{k} & =\sqrt{Ke^{i\phi}}=\sqrt{K}e^{i\phi/2}=\sqrt{\mu\omega\sqrt{(\epsilon\omega)^2+\sigma^2}}\left[\cos(\frac{1}{2}\arctan(\frac{\sigma}{\epsilon \omega}))+i\sin(\frac{1}{2}\arctan(\frac{\sigma}{\epsilon \omega}))\right] \\
           & =\omega\sqrt{\mu\epsilon}\sqrt{\sqrt{1+\left(\frac{\sigma}{\epsilon\omega}\right)^2}}\left[\cos(\frac{1}{2}\arctan(\frac{\sigma}{\epsilon \omega}))+i\sin(\frac{1}{2}\arctan(\frac{\sigma}{\epsilon \omega}))\right]
  \end{align*}
  and then use the half-angle identities to write this as
  $$
    \ti{k}=\omega\sqrt{\frac{\mu\epsilon}{2}}\sqrt{\sqrt{1+\left(\frac{\sigma}{\epsilon\omega}\right)^2}}\left[\sqrt{1+\cos(\arctan(\frac{\sigma}{\epsilon \omega}))}+i\sqrt{1-\cos(\arctan(\frac{\sigma}{\epsilon \omega}))}\right]
  $$
  and then apply the $\cos(\arctan(x))$ identity to rewrite that as
  $$
    \ti{k}=\omega\sqrt{\frac{\mu\epsilon}{2}}\sqrt{\sqrt{1+\left(\frac{\sigma}{\epsilon\omega}\right)^2}}\left[\sqrt{1+\frac{1}{\sqrt{1+\left(\frac{\sigma}{\epsilon \omega}\right)^2}}}+i\sqrt{1-\frac{1}{\sqrt{1+\left(\frac{\sigma}{\epsilon \omega}\right)^2}}}\right].
  $$
  Wow! Now multiply in the square root out front,
  $$
    \ti{k}=\omega\sqrt{\frac{\mu\epsilon}{2}}\left[\sqrt{\sqrt{1+\left(\frac{\sigma}{\epsilon\omega}\right)^2}+1}+i\sqrt{\sqrt{1+\left(\frac{\sigma}{\epsilon\omega}\right)^2}-1}\right].
  $$
  so
  $$\mathcal{k}=\omega\sqrt{\frac{\mu\epsilon}{2}}\sqrt{\sqrt{1+\left(\frac{\sigma}{\epsilon\omega}\right)^2}+1}$$
  and
  $$\kappa=\omega\sqrt{\frac{\mu\epsilon}{2}}\sqrt{\sqrt{1+\left(\frac{\sigma}{\epsilon\omega}\right)^2}-1}.$$
\end{soln}

% PROBLEM 2
\begin{problem}
Household microwave ovens operate at a frequency of
$\nu=\qty{2.54}{\giga\hertz}$. Imagine that
you leave a $\qty{20}{\degreeCelsius}$ stainless-steel spoon ($s=\qty{1.5e6}{\siemens\per\meter}$,
$\mu\approx\mu_0$, $\epsilon\approx\epsilon_0$) next to a
bowl of soup in the microwave.
\begin{enumerate}[label=(\alph*)]
  \item Estimate the proportion of the incident field that is reflected and
        transmitted from/into the spoon.
  \item Estimate the skin-depth over which the microwaves that enter the handle
        are absorbed
  \item Imagine that the full power of a $\qty{1}{\kilo\watt}$ microwave is incident onto the
        spoon's handle, after which any reflected power gets absorbed by the soup.
        If the spoon weighs $\qty{30}{\gram}$ and the steel has a heat capacity of $\qty{450}{\joule\per\kilo\gram\per\kelvin}$,
        what temperature do you expect the spoon to be at the end of a 2-minute
        full-power heating cycle?
\end{enumerate}
\end{problem}
\begin{soln}

\end{soln}
\end{document}